\section{Design Thesis}
\label{sec:thesis}

The design separates state reconciliation from state propagation. Each application packages an idempotent, commutative merge operation on byte-string state as a WebAssembly module. The platform determines where that state lives, how peers locate it, how replicas exchange differences, how updates reach subscribers, and how secrets remain local. Its routing and synchronization mechanisms operate independently of the application-supplied state algebra.

This separation distinguishes the system from two common approaches to decentralized state.

\paragraph{Application-defined rather than fixed CRDTs.} Conventional Conflict-free Replicated Data Type \citep{shapiro2011crdt} frameworks provide implementations of counters, sets, registers, sequences, and other predefined lattices. A Freenet contract supplies its own merge function, validity invariants, and tiebreak rules. This permits workload-specific state representations, but the runtime cannot statically verify their algebraic properties: a contract with a non-associative merge function will not converge. Because contracts are content-addressed, such a failure remains attributable to a specific module that users can reject.

\paragraph{Per-contract consistency instead of global consensus.} Blockchain platforms globally order state transitions. Freenet provides eventual consistency independently for each contract: replicas converge without coordinating the order of their updates with the rest of the network. This supports interactive applications whose updates can be merged, but not applications that require a single global order.

\paragraph{Separate public and private state.} A chat room's contents may be public while each user's signing key remains private. Contracts hold replicated public state; delegates hold device-local private state. Messages between them carry platform-verified sender identities across an enforced trust boundary.

The remainder of the paper describes the primitives that follow from this thesis (\cref{sec:primitives}), how updates move through the network (\cref{sec:updates}), how data is located (\cref{sec:routing}), how trust is established and what it does not cover (\cref{sec:trust}), the current implementation status and open problems (\cref{sec:status}), and the position in related work (\cref{sec:related}).
